\documentclass[11pt]{article}
\newcommand{\HomeworkHeader}{习题七}
% Shared LaTeX preamble for homework/*/main.tex

% --- Base layout ---
\usepackage[a4paper, top=2.5cm, bottom=2.5cm, left=2cm, right=2cm]{geometry}
\usepackage{fontspec}
\usepackage[UTF8]{ctex}%latex支持中文宏包
\usepackage{amsmath, amssymb, amsthm}
\usepackage{mathtools}
\usepackage[svgnames]{xcolor}
\usepackage{pgfplots}
\pgfplotsset{compat=1.18}
\usepackage[most]{tcolorbox}
\usepackage{enumitem}
% Modified: Change label from hyphen (-) to bullet point ($\bullet$) for better visibility
\setlist[itemize]{label=$\bullet$}
\setlist[enumerate]{label=(\arabic*)}


% --- Colors ---
\definecolor{primaryColor}{RGB}{46, 52, 64}
\definecolor{accentColor}{RGB}{94, 129, 172}
\definecolor{boxFill}{RGB}{236, 240, 241}
\definecolor{solLine}{RGB}{163, 190, 140}
\definecolor{expandColor}{RGB}{191, 97, 106}

% --- TikZ styles (DFA / NFA / PDA) ---
\usepackage{tikz}
\usetikzlibrary{automata, positioning, arrows.meta, bending, backgrounds}
\tikzset{
    dfa_state/.style={
        state,
        thick,
        draw=accentColor,
        fill=boxFill,
        text=primaryColor,
        minimum size=1.1cm
    },
    dfa_edge/.style={
        ->,
        >=stealth,
        thick,
        draw=primaryColor,
        auto
    },
    pda_node/.style={
        state,
        thick,
        draw=accentColor,
        fill=boxFill,
        text=primaryColor,
        minimum size=1.2cm
    },
    pda_edge/.style={
        ->,
        >=stealth,
        thick,
        draw=primaryColor,
        auto
    },
    rule_expand/.style={
        pda_edge,
        draw=expandColor,
        text=expandColor
    },
    rule_match/.style={
        pda_edge,
        draw=solLine,
        text=solLine!80!black
    }
}

% --- Header / footer ---
\usepackage{fancyhdr}
\pagestyle{fancy}
\setlength{\headheight}{13.6pt}%避免页眉高度警告
\fancyhf{}
\fancyhead[L]{\kaishu 中国科学院大学}
\fancyhead[C]{林俊杰2023K8009916012}
\fancyhead[R]{\kaishu 理论计算机}
\usepackage{lastpage}
\cfoot{\small 第 \thepage \ 页，共 \pageref{LastPage} 页}

% --- Problem box ---
\newtcolorbox[use counter=problemSeq]{problem}[1][]{
    enhanced,
    breakable,
    colback=boxFill,
    colframe=accentColor,
    coltitle=white,
    fonttitle=\bfseries\large,
    title={题目 ~\theproblemSeq \quad #1},
    boxrule=0.5mm,
    arc=3mm,
    drop shadow=black!15!white,
    attach boxed title to top left={xshift=0.5cm, yshift*=-3mm},
    boxed title style={colback=accentColor},
    % 关键修改：每当创建一个新问题盒子时，自动将步骤计数器重置为0
    code={\setcounter{stepcount}{0}}
}

% --- Solution 环境 (红色盒子样式，支持可选序号) ---
% 修改说明：
% 1. [1][] 表示接受一个可选参数，默认为空
% 2. title={...} 中使用 \ifstrempty 判断是否有参数，如果有则添加空格和参数
\newtcolorbox{solution}[1][]{
    enhanced, breakable,
    colback=red!3!white,        % 背景色：淡粉
    colframe=red!75!black,      % 边框色：深红
    coltitle=white,
    fonttitle=\bfseries\Large\itshape,
    title={Solution\ifstrempty{#1}{}{~#1}}, % 动态标题：Solution (1)
    boxrule=0.5mm,
    arc=3mm,
    drop shadow=black!15!white,
    attach boxed title to top left={xshift=0.5cm, yshift*=-3mm},
    boxed title style={colback=red!75!black},
    before skip=0.5em,          % 盒子前的间距
    after skip=2em,             % 盒子后的间距
    before upper={\setcounter{stepcount}{0}} % 每个解答开始时重置步骤计数
}

% --- Title ---
\newcommand{\makeCustomTitle}[1]{
    \begin{center}
        \vspace*{1cm}
        {\Huge \bfseries \color{primaryColor} 作业 #1}\\
        \vspace{0.5cm}
        {\Large \color{accentColor} \today}
        \vspace{1.2cm}
    \end{center}
}

% --- 自定义环境定义 ---

% --- 自定义计数器与环境 ---
\newcounter{stepcount}
\newcounter{casecount}[stepcount]
\newcounter{problemSeq} % 新增：专门用于 Problem 的独立计数器

% 【关键修改】挂载钩子：每当遇到 \section 命令（含 \section*），将 problemSeq 计数器重置为 0
\pretocmd{\section}{\setcounter{problemSeq}{0}}{}{}

% 2. 定义带自动编号的步骤环境
\newenvironment{proofstep}[1]
  {%
    \refstepcounter{stepcount}% 增加计数
    \par\vspace{1.5em}%
    \noindent\textbf{\large \thestepcount. #1}% 输出 "1. 标题"
    \par\vspace{0.5em}%
  }
  {\par}

% 3. 定义无编号的结论环境 (用于最后的结论，格式同步骤)
\newenvironment{proofresult}[1]
  {%
    \par\vspace{1.5em}%
    \noindent\textbf{\large #1}% 仅输出标题
    \par\vspace{0.5em}%
  }
  {\par}

% 4. 定义带自动编号的情形环境
\newenvironment{proofcase}[1]
  {%
    \refstepcounter{casecount}% 增加计数
    \par\vspace{1em}%
    \noindent\textbf{情形 \thecasecount：#1}% 输出 "情形 1：标题"
    \par\vspace{0.2em}%
  }
  {\par}

\begin{document}
\makeCustomTitle{习题七}
% 来源：assets/习题七.pdf，红框原题 2、5。
% 原题 2 跨 PDF 第 1、2 页（教材第 164、165 页）；
% 原题 5 在 PDF 第 2 页（教材第 165 页）。第 3 页无红框题。
% 两道选中题目均无附图。

\begin{problem}[原题 2]
给定系统和初始条件
\[
 \dot x=\begin{bmatrix}0&1\\0&0\end{bmatrix}x
 +\begin{bmatrix}0\\1\end{bmatrix}u,\qquad
 x(0)=\begin{bmatrix}1\\2\end{bmatrix},
\]
以及性能指标
\[
 J=\int_0^\infty(2x_1^2+2x_1x_2+x_2^2+u^2)\,dt.
\]
求最优控制及最优性能指标值。
\end{problem}
\begin{solution}
\begin{proofstep}{识别性能指标权重并检查适用条件}
将性能指标写成 $J=\int_0^\infty(x^{\mathsf T}Qx+u^{\mathsf T}Ru)\,dt$，其中
\[
 Q=\begin{bmatrix}2&1\\1&1\end{bmatrix}>0,\qquad R=1.
\]
交叉项 $2x_1x_2$ 对应 $Q_{12}=Q_{21}=1$。
这是因为 $x^{\mathsf T}Qx=Q_{11}x_1^2+2Q_{12}x_1x_2+Q_{22}x_2^2$。
由 $Q_{11}=2>0$、$\det Q=2\times1-1=1>0$，可知 $Q$ 正定。
$[B\;AB]=\begin{bmatrix}0&1\\1&0\end{bmatrix}$ 满秩，
系统可控，且 $Q>0$ 保证可检测，因此存在唯一正定的稳定化 Riccati 解。
原题积分前没有 $\tfrac12$，下文取值函数 $V(x)=x^{\mathsf T}Px$。
\end{proofstep}

\begin{proofstep}{建立 Riccati 方程并逐项展开}
因为 $\nabla V=2Px$，Hamilton--Jacobi--Bellman 方程中的控制项为
\[
 u^{\mathsf T}Ru+2x^{\mathsf T}PBu.
\]
对 $u$ 求导，$2Ru+2B^{\mathsf T}Px=0$，从而
$u^*=-R^{-1}B^{\mathsf T}Px$。
将它代回值函数方程，得到代数 Riccati 方程
\[
 A^{\mathsf T}P+PA-PBR^{-1}B^{\mathsf T}P+Q=0
\]
令 $P=\begin{bmatrix}p&q\\q&r\end{bmatrix}$，本题各矩阵乘积为
\[
 A^{\mathsf T}P+PA=\begin{bmatrix}0&p\\p&2q\end{bmatrix},\qquad
 PBR^{-1}B^{\mathsf T}P=\begin{bmatrix}q^2&qr\\qr&r^2\end{bmatrix}.
\]
故方程的 $(1,1),(1,2),(2,2)$ 三个独立分量分别给出
\[
 \begin{cases}
 2-q^2=0,\\
 p-qr+1=0,\\
 2q-r^2+1=0.
 \end{cases}
\]
\end{proofstep}
\begin{proofstep}{选择稳定化分支并验证正定性}
由第一式可得 $q=\pm\sqrt2$；再由第三式得 $r^2=1+2q$，
第二式给出 $p=qr-1$。
选择根号的符号时，需要检查最优闭环的稳定性，而非任取代数根。
本题
\[
 A-BB^{\mathsf T}P=\begin{bmatrix}0&1\\-q&-r\end{bmatrix},\qquad
 \det(sI-A+BB^{\mathsf T}P)=s^2+rs+q.
\]
二阶 Hurwitz 条件是 $r>0,q>0$，故必须取
\[
 q=\sqrt2,\qquad r=\sqrt{1+2\sqrt2},\qquad p=qr-1.
\]
此外 $q=-\sqrt2$ 会使 $r^2=1-2\sqrt2<0$，无法给出实对称解。
为简写，令 $a=\sqrt2$、$b=\sqrt{1+2\sqrt2}$。
于是
\[
 \boxed{P=\begin{bmatrix}ab-1&a\\a&b\end{bmatrix},\qquad
 u^*(t)=-a x_1(t)-b x_2(t).}
\]
按教材 $u=Kx$ 的符号约定，$K=[-a\;-b]$。
此时 $(ab)^2=2(1+2\sqrt2)>1$，故 $p=ab-1>0$。
第二个顺序主子式为
\[
 \det P=(ab-1)b-a^2=ab^2-b-2
 =\sqrt2(1+2\sqrt2)-b-2=2+\sqrt2-b.
\]
又因 $(2+\sqrt2)^2-b^2=5+2\sqrt2>0$ 且二者均正，
所以 $\det P>0$，从而 $P>0$。
由二次方程求根公式得到闭环极点
\[
 \boxed{\lambda_{1,2}=-\frac b2\pm\frac{j}{2}\sqrt{2\sqrt2-1},}
\]
实部严格为负，验证了所取的是稳定化解。
\end{proofstep}

\tcbbreak
\begin{proofstep}{证明最优性并代入初值计算最优代价}
令 $V=x^{\mathsf T}Px$，沿原系统有
\[
 \dot V=x^{\mathsf T}(A^{\mathsf T}P+PA)x
 +2uB^{\mathsf T}Px.
\]
利用 Riccati 方程代换 $A^{\mathsf T}P+PA=PBB^{\mathsf T}P-Q$，整理为
\[
 x^{\mathsf T}Qx+u^2
 =-\dot V+(u+B^{\mathsf T}Px)^2.
\]
对于趋于原点的可容许轨迹，对上式从 $0$ 到 $\infty$ 积分：
\[
 J=x(0)^{\mathsf T}Px(0)
 +\int_0^\infty (u+B^{\mathsf T}Px)^2\,dt
 \geq x(0)^{\mathsf T}Px(0).
\]
取 $u=u^*=-B^{\mathsf T}Px$ 时，平方项恒为零，且已证闭环稳定，
故下界可达到。这也说明本题的最优值没有额外系数 $\tfrac12$。
将 $x(0)=[1\;2]^{\mathsf T}$ 代入：
\[
 \begin{aligned}
 J^*&=\begin{bmatrix}1&2\end{bmatrix}
       \begin{bmatrix}p&q\\q&r\end{bmatrix}\begin{bmatrix}1\\2\end{bmatrix}\\
 &=p+4q+4r=(ab-1)+4a+4b=(a+4)b+4a-1.
 \end{aligned}
\]
因此
\[
 \boxed{J^*=x(0)^{\mathsf T}Px(0)
 =(4+\sqrt2)\sqrt{1+2\sqrt2}+4\sqrt2-1.}
\]
\end{proofstep}
\begin{proofstep}{写出给定初值对应的最优状态轨迹}
最优闭环满足
\[
 \dot x_1^*=x_2^*,\qquad \dot x_2^*=-ax_1^*-bx_2^*,
\]
故 $\ddot x_1^*+b\dot x_1^*+ax_1^*=0$。
记 $\omega=\tfrac12\sqrt{2\sqrt2-1}$，通解为
$x_1^*=e^{-bt/2}(c_1\cos\omega t+c_2\sin\omega t)$。
由 $x_1^*(0)=1$、$\dot x_1^*(0)=x_2^*(0)=2$，有
\[
 c_1=1,\qquad -\frac b2c_1+\omega c_2=2,
 \qquad c_2=\frac{2+b/2}{\omega}.
\]
所以
\[
 x_1^*(t)=e^{-bt/2}\left[
 \cos(\omega t)+\frac{2+b/2}{\omega}\sin(\omega t)\right],
 \qquad x_2^*(t)=\dot x_1^*(t),
\]
相应最优控制为 $u^*(t)=-ax_1^*(t)-bx_2^*(t)$。
\end{proofstep}
\end{solution}

\clearpage
\begin{problem}[原题 5]
给定系统
\[
 \dot x=\begin{bmatrix}1&0\\0&2\end{bmatrix}x
 +\begin{bmatrix}1\\1\end{bmatrix}u,\qquad
 x(0)=\begin{bmatrix}2\\1\end{bmatrix},\qquad
 y=\begin{bmatrix}1&2\end{bmatrix}x,
\]
以及性能指标
\[
 J=\int_0^\infty(y^2+2u^2)\,dt.
\]
求最优控制、最优性能指标值及最优闭环系统的特征值。
\end{problem}
\begin{solution}
\begin{proofstep}{由输出代价确定状态权重并检验条件}
因为 $y=x_1+2x_2$，有 $y^2=x_1^2+4x_1x_2+4x_2^2$，故
\[
 Q=C^{\mathsf T}C=\begin{bmatrix}1&2\\2&4\end{bmatrix},\qquad R=2.
\]
虽然 $Q$ 仅半正定，但
\[
 \det[B\;AB]=\det\begin{bmatrix}1&1\\1&2\end{bmatrix}=1,
 \qquad
 \det\begin{bmatrix}C\\CA\end{bmatrix}
 =\det\begin{bmatrix}1&2\\1&4\end{bmatrix}=2,
\]
所以系统可控且可观测，满足无限时域 LQR 的稳定化条件。
这里 $R=2$ 是输入二次项的系数。
原题没有 $\tfrac12$，采用 $V=x^{\mathsf T}Px$ 时最优反馈仍为
$u^*=-R^{-1}B^{\mathsf T}Px=-\tfrac12 B^{\mathsf T}Px$。
\end{proofstep}

\begin{proofstep}{展开代数 Riccati 方程}
令
\[
 P=\begin{bmatrix}p&q\\q&r\end{bmatrix},\qquad
 F=R^{-1}B^{\mathsf T}P=\begin{bmatrix}f_1&f_2\end{bmatrix},
 \quad f_1=\frac{p+q}{2},\quad f_2=\frac{q+r}{2}.
\]
先计算
\[
 A^{\mathsf T}P+PA=\begin{bmatrix}2p&3q\\3q&4r\end{bmatrix},\qquad
 PB=\begin{bmatrix}p+q\\q+r\end{bmatrix},
\]
\[
 \frac12PBB^{\mathsf T}P
 =\frac12\begin{bmatrix}(p+q)^2&(p+q)(q+r)\\(p+q)(q+r)&(q+r)^2\end{bmatrix}.
\]
因此 Riccati 方程
$A^{\mathsf T}P+PA-\tfrac12PBB^{\mathsf T}P+Q=0$ 的各独立分量为
\[
 \begin{cases}
 2p-\tfrac12(p+q)^2+1=0,\\
 3q-\tfrac12(p+q)(q+r)+2=0,\\
 4r-\tfrac12(q+r)^2+4=0.
 \end{cases}
\]
利用 $p+q=2f_1,q+r=2f_2$，可进一步写成
\[
 2p-2f_1^2+1=0,\qquad
 3q-2f_1f_2+2=0,\qquad
 4r-2f_2^2+4=0.
\]
这些代数方程可能有多组解，下面用闭环稳定性选出所需的一组。
\end{proofstep}
\begin{proofstep}{计算 Hamilton 矩阵的特征多项式}
考虑 Hamilton 矩阵
\[
 \mathcal H=\begin{bmatrix}A&-BR^{-1}B^{\mathsf T}\\-Q&-A^{\mathsf T}\end{bmatrix}
 =\begin{bmatrix}
 1&0&-\frac12&-\frac12\\
 0&2&-\frac12&-\frac12\\
 -1&-2&-1&0\\
 -2&-4&0&-2
 \end{bmatrix}.
\]
下面给出行列式的计算过程。记
\[
 D=sI-A=\operatorname{diag}(s-1,s-2),\qquad
 E=sI+A^{\mathsf T}=\operatorname{diag}(s+1,s+2).
\]
则
\[
 sI-\mathcal H=\begin{bmatrix}D&\tfrac12BB^{\mathsf T}\\C^{\mathsf T}C&E\end{bmatrix}.
\]
先对 $D$ 作 Schur 补，再对秩一修正使用行列式公式，得到
\[
 \begin{aligned}
 \det(sI-\mathcal H)
 &=\det D\,\det\left(E-\tfrac12C^{\mathsf T}CD^{-1}BB^{\mathsf T}\right)\\
 &=\det D\,\det E\left[1-\tfrac12(CD^{-1}B)(B^{\mathsf T}E^{-1}C^{\mathsf T})\right].
 \end{aligned}
\]
其中两个标量为
\[
 CD^{-1}B=\frac1{s-1}+\frac2{s-2}
 =\frac{3s-4}{(s-1)(s-2)},
\]
\[
 B^{\mathsf T}E^{-1}C^{\mathsf T}=\frac1{s+1}+\frac2{s+2}
 =\frac{3s+4}{(s+1)(s+2)}.
\]
代入并约去分母，得到多项式恒等式
\[
 \begin{aligned}
 \det(sI-\mathcal H)
 &=(s^2-1)(s^2-4)-\frac12(3s-4)(3s+4)\\
 &=s^4-5s^2+4-\frac92s^2+8\\
 &=s^4-\frac{19}{2}s^2+12.
 \end{aligned}
\]
以上计算最初在 $s\ne\pm1,\pm2$ 时使用逆矩阵；化成多项式后恒等式对所有 $s$ 成立。
令 $z=s^2$，则 $z^2-\tfrac{19}2z+12=0$，即 $2z^2-19z+24=0$。
按二次方程求根公式，两根为
\[
 z=\frac{19\pm\sqrt{19^2-4\times24\times2}}4
 =\frac{19\pm13}{4}=8,\ \frac32.
\]
因此
\[
 \det(sI-\mathcal H)=(s^2-8)\left(s^2-\frac32\right).
\]
\end{proofstep}
\begin{proofstep}{选择稳定特征值，并由系数匹配求反馈增益}
Riccati 方程可以重排为
$-Q-A^{\mathsf T}P=P(A-BR^{-1}B^{\mathsf T}P)$，因而
\[
 \mathcal H\begin{bmatrix}I\\P\end{bmatrix}
 =\begin{bmatrix}I\\P\end{bmatrix}(A-BF),
\]
该式说明 $\begin{bmatrix}I\\P\end{bmatrix}$ 的列空间是
$\mathcal H$ 的不变子空间，$A-BF$ 的两个特征值来自 $\mathcal H$。
由于要求最优闭环渐近稳定，必须取两枚负特征值；取正根的代数分支不符合稳定化要求。
记
\[
 a=2\sqrt2,\qquad b=\frac{\sqrt6}{2},\qquad
 \sigma=a+b,\qquad \rho=ab=2\sqrt3.
\]
则所需闭环多项式为 $(s+a)(s+b)=s^2+\sigma s+\rho$。
另一方面
\[
 A-BF=\begin{bmatrix}1-f_1&-f_2\\-f_1&2-f_2\end{bmatrix},
\]
\[
 \det(sI-A+BF)=s^2+(f_1+f_2-3)s+2-2f_1-f_2.
\]
因此比较 $s$ 项和常数项系数得
\[
 f_1+f_2=3+\sigma,\qquad 2f_1+f_2=2-\rho.
\]
第二式减第一式，得到 $f_1=-1-\sigma-\rho$；
再代回第一式求得 $f_2=4+2\sigma+\rho$，即
\[
 \boxed{f_1=-1-\sigma-\rho,\qquad f_2=4+2\sigma+\rho.}
\]
\end{proofstep}
\begin{proofstep}{恢复 Riccati 解并给出最优反馈}
由第 2 步的三个 Riccati 分量方程，依次求得
\[
 p=f_1^2-\frac12,\qquad
 q=\frac{2f_1f_2-2}{3},\qquad
 r=\frac{f_2^2}{2}-1.
\]
展开为
\[
 \begin{aligned}
 p&=22+10\sqrt2+8\sqrt3+9\sqrt6,\\
 q&=-24-14\sqrt2-12\sqrt3-10\sqrt6,\\
 r&=32+22\sqrt2+16\sqrt3+12\sqrt6.
 \end{aligned}
\]
需要同时检查恢复出的 $P$ 与原增益定义一致。直接相加可得
\[
 \begin{aligned}
 p+q&=-2-4\sqrt2-4\sqrt3-\sqrt6
     =2(-1-\sigma-\rho)=2f_1,\\
 q+r&=8+8\sqrt2+4\sqrt3+2\sqrt6
     =2(4+2\sigma+\rho)=2f_2.
 \end{aligned}
\]
所以 $F=\tfrac12 B^{\mathsf T}P$ 的定义成立，三个分量方程也随之同时满足。
此外 $p>0$ 且
\[
 \det P=176+132\sqrt2+108\sqrt3+72\sqrt6>0,
\]
故 $P>0$。所求最优控制为
\[
 \boxed{u^*(t)=(1+\sigma+\rho)x_1(t)
 -(4+2\sigma+\rho)x_2(t).}
\]
按教材的 $u=Kx$ 约定，反馈矩阵为
$K=-F=\begin{bmatrix}1+\sigma+\rho&-4-2\sigma-\rho\end{bmatrix}$。
\end{proofstep}

\begin{proofstep}{完成平方，求最优性能指标并核对闭环特征值}
令 $V=x^{\mathsf T}Px$，利用 Riccati 方程和 $B^{\mathsf T}P=2F$，有
\[
 \begin{aligned}
 \dot V
 &=x^{\mathsf T}(A^{\mathsf T}P+PA)x+2x^{\mathsf T}PBu\\
 &=x^{\mathsf T}(2F^{\mathsf T}F-Q)x+4(Fx)u\\
 &=2(Fx)^2-y^2+4(Fx)u.
 \end{aligned}
\]
整理后可写成
\[
 y^2+2u^2=-\dot V+2(u+Fx)^2.
\]
沿趋于原点的可容许轨迹积分，得
\[
 J=x(0)^{\mathsf T}Px(0)+2\int_0^\infty(u+Fx)^2\,dt.
\]
取 $u=-Fx$ 时，平方项为零，稳定闭环保证 $V(\infty)=0$，故达到最小值。
代入 $x(0)=[2\;1]^{\mathsf T}$，先计算
\[
 J^*=\begin{bmatrix}2&1\end{bmatrix}
 \begin{bmatrix}p&q\\q&r\end{bmatrix}\begin{bmatrix}2\\1\end{bmatrix}
 =2(2p+q)+(2q+r)=4p+4q+r.
\]
再把各根式的系数分别相加：
\[
 \begin{aligned}
 4p+4q+r
 &=(88-96+32)+(40-56+22)\sqrt2\\
 &\quad +(32-48+16)\sqrt3+(36-40+12)\sqrt6\\
 &=24+6\sqrt2+8\sqrt6.
 \end{aligned}
\]
所以
\[
 \boxed{J^*=\begin{bmatrix}2&1\end{bmatrix}P
 \begin{bmatrix}2\\1\end{bmatrix}=4p+4q+r
 =24+6\sqrt2+8\sqrt6\approx52.0812.}
\]
最优闭环的特征值为
\[
 \boxed{\lambda_1=-2\sqrt2,\qquad
 \lambda_2=-\frac{\sqrt6}{2}.}
\]
两者均严格为负，与选择的 Hamilton 稳定子空间一致。
\end{proofstep}
\tcbbreak
\begin{proofstep}{求给定初值对应的最优控制时间函数}
由于闭环有两个不同的实极点 $-a,-b$，最优控制具有形式
\[
 u^*(t)=c_a e^{-at}+c_b e^{-bt}.
\]
利用已求出的反馈增益，先计算初值
\[
 u^*(0)=-F\begin{bmatrix}2\\1\end{bmatrix}
 =-(2f_1+f_2)=\rho-2.
\]
又有
\[
 \begin{aligned}
 \dot u^*(0)&=-F(A-BF)x(0)\\
 &=-2(f_1+f_2)+(f_1+f_2)(2f_1+f_2)\\
 &=-2(3+\sigma)+(3+\sigma)(2-\rho)
 =-\rho(3+\sigma).
 \end{aligned}
\]
因此 $c_a,c_b$ 满足
\[
 c_a+c_b=\rho-2,\qquad ac_a+bc_b=\rho(3+\sigma).
\]
第一式乘以 $b$ 并从第二式减去，得到
\[
 c_a=\frac{\rho(3+\sigma)-b(\rho-2)}{a-b}
 =\frac{\rho(3+a)+2b}{a-b}.
\]
再用 $c_b=\rho-2-c_a$，得到
$c_b=-[\rho(3+b)+2a]/(a-b)$。所以
\[
 \boxed{u^*(t)=
 \frac{\rho(3+a)+2b}{a-b}e^{-at}
 -\frac{\rho(3+b)+2a}{a-b}e^{-bt}.}
\]
该时间函数与状态反馈 $u^*=-Fx^*$ 等价，并随时间趋于零。
\end{proofstep}
\end{solution}
\end{document}
